\documentclass[12pt,a4paper]{article}
%per compilare il template usare questo sotto e commentare l'altro:
%\usepackage{../stile}
\usepackage{../../../template/stile}

%Titolo documento
\newcommand{\titoloDocumento}{Template documenti per SWE}

%Prima data di creazione del documento
\newcommand{\dataCreazione}{30 Novembre 2015}

%Inserite la versione attuale del documento
\newcommand{\versione}{1.0.7}

%Stato in cui si trova il documento: Formale solo all'atto di consegna
\newcommand{\stato}{Informale}

%Uso del documento
\newcommand{\uso}{Interno}

\rhead{\titoloDocumento}
\lfoot{Versione: \versione}
\title{\titoloDocumento}

\begin{document}
\begin{titlepage}
\begin{center}
\today \\
\vspace{1cm}
\begin{Huge}
\textbf{\nomeGruppo} \\
\end{Huge}
\textbf{\prjL} \\
\vspace{1cm}
\includegraphics[scale=0.5]{\logoGrande}
\vspace{1cm}

\HRule \\[0.4cm]
\begin{Huge}
{\huge \bfseries \titoloDocumento}\\[0.4cm]
\end{Huge}
\HRule \\[1cm]
\vfill

\begin{table}[h]
\begin{center}
\begin{tabular}{r | l}
\multicolumn{2}{c}{\textbf{Informazioni sul documento}}\\
\midrule
\textbf{Nome Documento}	&	\titoloDocumento	\\
\textbf{Versione}	&	\versione	\\
\textbf{Stato}	&	\emph{\stato}	\\
\textbf{Uso}	&	\emph{\uso}	\\
\textbf{Data Creazione}	&	\dataCreazione	\\
\textbf{Data Ultima Modifica}	&	\today	\\
\textbf{Redazione}	&	\IB\\
\ 	&	\TP\\
\ 	&		\AVE\\
\textbf{Verifica}	&		\NDC\\
\ & \AVI \\
\textbf{Approvazione}	&		\WS\\
\textbf{Lista Distribuzione}	&	\nomeGruppo	\\
\ 	&	\Vardanega	\\
\ 	&	\Cardin	\\
\ 	&	Il proponente \Zucchetti	\\

\end{tabular}
\end{center}
\end{table}

\end{center}
\end{titlepage}
\newpage

\Large{\textbf{Registro delle modifiche}}\\
\normalsize

\begin{table}[h]
\begin{center}

\begin{tabular}{p{0.12\textwidth} p{0.25\textwidth} p{0.18\textwidth} p{0.4\textwidth}}
\toprule
\textbf{Versione}	&	\textbf{Autore}	&	\textbf{Data}	&	\textbf{Descrizione}\\
\midrule
\midrule
v1.0.0	&	\NDC	&	30 Novembre 2015 	&	Prima revisione\\
\midrule
v1.0.1	&	\TP	&	6 Dicembre 2015	&	Aggiunto capitolato 1\\
\midrule
v1.0.2	&	\TP	&	7 Dicembre 2015	&	Aggiunto capitolato 5\\
\midrule
v1.0.3	&	\IB	&	8 Dicembre 2015	&	Aggiunto capitolato 3\\
\midrule
v1.0.4	&	\IB	&	8 Dicembre 2015	&	Aggiunto capitolato 4\\
\midrule
v1.0.5	&	\AVE	&	8 Dicembre 2015	&	Aggiunto capitolato 2 e 6\\
\midrule
v1.0.6  & \TP & 18 Dicembre 2015 & Normalizzate sezioni capitolati\\
\midrule
v1.0.7  & \TP & 18 Dicembre 2015 & Aggiunta introduzione\\
\bottomrule
\end{tabular}
\caption{Versionamento del documento}
\label{tabVers1}
\end{center}
\end{table}
\newpage

\tableofcontents
\newpage

\section{Introduzione}
\subsection{Scopo del Documento}

In questo documento vengono analizzati i requisiti, aspetti positivi e fattori di rischio
di ogni capitolato per evidenziare i punti di forza e gli aspetti negativi che hanno portato ad accettare
o scartare rispettivamente un capitolato.

\subsection{Riferimenti}
\subsubsection{Informativi}

\begin{itemize}
\item \href{http://www.math.unipd.it/~tullio/IS-1/2015/Progetto/C1.pdf}{Capitolato C1}
\item \href{http://www.math.unipd.it/~tullio/IS-1/2015/Progetto/C2.pdf}{Capitolato C2}
\item \href{http://www.math.unipd.it/~tullio/IS-1/2015/Progetto/C3.pdf}{Capitolato C3}
\item \href{http://www.math.unipd.it/~tullio/IS-1/2015/Progetto/C4.pdf}{Capitolato C4}
\item \href{http://www.math.unipd.it/~tullio/IS-1/2015/Progetto/C5.pdf}{Capitolato C5}
\item \href{http://www.math.unipd.it/~tullio/IS-1/2015/Progetto/C6.pdf}{Capitolato C6}
\end{itemize}


\newpage
\section{Capitolato C5 - Quizzpedia}
\subsection{Descrizione}

Il capitolato propone la creazione di un sito web per la creazione e somministrazione di questionari scelti
per argomento.

\subsubsection{Dominio applicativo}

Le domande verranno fornite attraverso uno specifico \textit{Quiz Markup Language} (QML), la cui creazine
è oggetto principale del capitolato.
L'interfaccia amministrativa dovrà supportare dispositivi desktop, mentre per i candidati viene richiesto 
il supporto di dispositivi desktop e mobili.

Si dovranno considerare tre tipi di profili:

\begin{itemize}
\item Amministratore (crea profili, gestisce questionari e domande);
\item Insegnante (crea domande e questionari e analizza i risultati);
\item Studente (completa i questionari).
\end{itemize}

con i relativi permessi.

\subsubsection{Dominio tecnologico}

Per quanto riguarda il server è richiesto l'utilizzo di una delle seguenti opzioni:

\begin{itemize}
\item Tomcat (Java);
\item Node.js (JavaScript).
\end{itemize}

Per l'archiviazione sono proposte le seguenti tecnologie:

\begin{itemize}
\item PostgreSQL;
\item MongoDB.
\end{itemize}

Per il client viene richiesto l'utilizzo di HTML5, CSS3 e il supporto desktop per l'intera applicazione e mobile
solo per la parte destinata allo studente.

Per la creazione del linguaggio QML il gruppo ha considerato l'utilizzo di XML o altri linguaggi di markup
estendibili.

\subsection{Valutazione}
\subsubsection{Aspetti positivi}

Sono stati evidenziati i seguenti aspetti positivi che hanno contribuito alla scelta di
questo capitolato:

\begin{itemize}
\item semplicità e chiarezza degli obiettivi;
\item la gestione dell'interfaccia \textit{responsive} può essere risolto attraverso i numerosi
    framework (tra cui \href{http://getbootstrap.com}{Bootstrap} e \href{http://foundation.zurb.com}{Foundation});
\item PostgreSQL e il sistema relazionale è propedeutico per questo corso quindi è conosciuto dall'intero gruppo;
\item HTML, CSS e Javascript sono trattate nel corso di Tecnologie Web che tutti i membri del gruppo stanno o hanno
seguito;
\item il cuore applicativo è semplice da implementare ed è possibile utilizzare un approccio incrementale
per lo sviluppo delle varie funzioni, sia richieste che opzionali.
\end{itemize}

\subsubsection{Potenziali criticità}

Sono state inoltre rilevate possibili punti critici nello sviluppo applicativo e/o nelle tecnologie utilizzate
che dovranno essere tenuti sotto controllo:

\begin{itemize}
\item nel caso di utilizzo di MongoDB, la tecnologia il paradigma (NoSQL) non sono affatto conosciuti;
\item utilizzando Javascript bisogna tenere a mente che la semplicità e la dinamicità del linguaggio
sono armi a doppio taglio;
\item sia Tomcat che Node.js sono sconosciuti a tutti i membri del gruppo;
\item le tecnologie utilizzate sono per lo più già conosciute e non contribuiscono molto
al curriculum individuale.
\end{itemize}

Per lo sviluppo Javascript, vengono consigliati strumenti per un maggiore controllo della correttezza del
codice (come \href{http://flowtype.org}{Flow}) oppure altri linguaggi tipati che compilati producono Javascript
(come \href{http://coffeescript.org}{CoffeeScript}).

\subsection{Conclusione}

Nonostante alcune tecnologie siano sconosciute o presentino aspetti negativi, il capitolato è abbastanza flessibile
da permettere di limitare tali rischi scegliendo tecnologie più sicure o strumenti che rendano più difficile commettere
errori.

Il dominio applicativo risulta chiaro e più in linea con il tipo di progetti a cui siamo stati sottoposti da altri corsi. 
Per cui crediamo che gli obiettivi e le tecnologie utilizzate
ci diano una maggiore sicurezza, al costo dell'esposizione a tecnologie e modelli meno stimolanti.

\newpage
\section{Altri Capitolati}

\subsection{Capitolato C1 Actorbase}
\subsubsection{Descrizione}

Il capitolato richiede l'implementazione di un \textit{database NoSQL} di tipo \textit{key-value} basato
sull'architettura ad attori. E' inoltre richiesta la definizione di un \textit{domain-specific-language} (DSL)
per poter interagire con il \textit{database} da linea da comando.

\subsubsection{Aspetti positivi}

Il gruppo ha rilevato i seguenti aspetti positivi:

\begin{itemize}
\item Scala risulta un linguaggio interessante e conciso;
\item l'architettura proposta non è di difficile comprensione;
\item il modello \textit{actor-model} è presentato anche da Programmazione Concorrente e Distribuita di quest'anno, corso che è seguito dalla maggiorparte del gruppo;
\item possibilità di scegliere Java al posto di Scala, considerando che tutto il gruppo ha asperienza con il primo;
\item il proponente interno potrebbe avere una maggiore disponibilità.
\end{itemize}

\subsubsection{Fattori di rischio}

Questo capitolato non è stato scelto per i seguenti motivi:

\begin{itemize}
\item la libreria Akka è per lo più sconosciuta ai membri del gruppo;
\item nel caso si scelga Java, l'interazione con la libreria Akka risulterebbe di gran lunga più verbosa
in quanto essa è stata pensata per Scala;
\item poca conoscenza delle caratteristiche dei sistemi distribuiti;
\item lo sviluppo di un DSL appropriato potrebbe richiedere conoscenze relative al \textit{parsing} e 
analisi lessicale;
\item se Scala è scelto come linguaggio implementativo, sarà necessario imparare il linguaggio e l'ecosistema
di strumenti e librerie a sua dispozione.
\end{itemize}

\newpage

\subsection{Capitolato C2 CLIPS}

\subsubsection{Descrizione}

Per questo capitolato è richiesta l'implementazione di un'applicazione per Smartphone. L'applicazione deve permettere di gestire la microlocalizzazione di luoghi fisici, dove è presente la tecnologia BLE Beacon, e veicolare al utente finale informazioni contestuali a questi luoghi. 

\subsubsection{Aspetti positivi}

Il gruppo ha trovato questo capitolato indubbiamente interessante perchè  posto nell'ambito della ricerca,permettendo lo studio di tecnologie recenti, e anche perchè viene data la possibilità di proporre e sviluppare una propria idea originale.

\subsubsection{Fattori di rischio}

Tuttavia a causa una serie di limiti legati alla tecnologia Beacon BLE, il team ha deciso di propendere verso un altro capitolato per i seguenti motivi:

\begin{itemize}
\item La teconologia dei Beacon BLE e le sue applicazioni sono ancora immature;
\item Abbiamo ritenuto che un costo medio di 25 euro per Beacon fosse eccessivamente mettendolo in relazione con le funzionalità offerte;
\item L'obbligo da parte del sistema operativo Android o IOS di tenere attivo il sistema di localizzazione dello Smartphone per poter localizzare i Beacon, risultando in un consumo notevole della batteria;
\item L'impossibilità di interagire in alcun modo con i Beacon senza un applicazione dedicata.
\item Difficoltà nel proporre un idea non banale e funzionale considerando l'imprecisione e il tempo di latenza nel rilevamento dei Beacon da parte si uno smartphone.
\end{itemize}

\newpage
\subsection{Capitolato C3 UMAP}

\subsubsection{Descrizione}

Questo capitolato richiede l'implementazione di un sistema per la  raccolta e la elaborazione di dati eterogenei, dando inoltre la possibilità all'utente finale di leggere I dati e l'analisi fatta sugli stessi da un engine predittivo.

Il capitolato richiede lo sviluppo di:
\begin{itemize}
\item Frontend;
\item Business Logic API; 
\item Engine predittivo;
\item Broker MQTT. 
\end{itemize}

\subsubsection{Aspetti positivi}

E' stato rilevato solo il seguente aspetto positivo:

\begin{itemize}
\item	Le tecnologie necessarie per svolgere il capitolato sono svariate (java, mongoDB, AWS , html e molte altre ) il che permetterebbe al gruppo di venire in contatto con strumenti nuovi ed estremamente utili in ambito lavorativo.
\end{itemize}

\subsubsection{Fattori di rischio}

\begin{itemize}
\item	Lo studio e l'apprendimento delle tecnologie necessarie sicuramente risulterebbe in uno sforzo notevole da parte di ogni componente del gruppo;

\item	Un altra difficolta risiede nella realizzazione del motore predittivo dato che al gruppo mancano totalmente le idee e le conoscenze per concepire l'elaborazione di dati eterogenei, considerando anche il fatto che durante la presentazione non sono	stati dati spunti a riguardo;
\item	Il sistema da realizzare sarebbe incredibilmente amplio.
\end{itemize}

\newpage
\subsection{Capitolato C4 MaaS}
\subsubsection{Descrizione}
Maas è basato su un prodotto di nome Maap, che è un database noSql per il trattamento di dati Business attraverso un interfaccia grafica permettendone l'utilizzo a persone che non sono esperte di database.Il committente richiede che questo prodotto diventi un servizio gestito attraverso un account web che dia la possibilità di eliminare l'onere all'utente di eseguire la preparazione e l'installazione di Maap.

Sono stati rilevati i seguenti requisiti:

\begin{itemize}

\item	Le tecnologie richieste includono:

\begin{itemize}
\item Node.js per il backend. MaaS deve essere in Long Term Support [LTS] versione Argon;
\item	MongoDB, versione maggiore di 3.x come database per l'applicazione; 
\end{itemize}

\item Il Sistema deve essere costruito usando un framework di alto livello per Node.js. Noi fortemente consigliamo di usare IBM loopback;

\item Il sistema deve essere schierabile con Heroku;
 
\item	Il codice sorgente deve essere pubblicato e versionato   usando github oppure bitbucket. 

\end{itemize}

\subsubsection{Aspetti positivi}
\begin{itemize}
\item	Potenzialmente la mole di lavoro e la complessità dello stesso potrebbero essere esigue dato che il capitolato si basa già su un prodotto esistente;

\item	Alcune delle tecnologie a supporto del capitolato potrebbero essere , anche se in modo parziale, già note al gruppo ( html, css, un servizio di	database, java ).
\end{itemize}

\subsubsection{Fattori di rischio}
\begin{itemize}
\item	L'argomento risulta essere poco stimolante, dato che la parte più interessante sarebbe stata la creazione di Maap che è già stato implementato;

\item	Essendo il committente una Startup con risorse di tempo, personale e	finanze limitate è possibile che sia difficile aiutare e seguire in	gruppo di lavoro in caso di problemi o dubbi inerenti al progetto.
\end{itemize}


\newpage
\subsection{Capitolato C6 MIVOQ}
\subsubsection{Descrizione}

Il Capitolato richiede che il gruppo di lavoro proponga un idea e realizzi applicazione mobile che sfrutti le funzionalità di text to speech fornite attraverso delle API HTTP.

\subsubsection{Aspetti positivi}
\begin{itemize}
	\item Le tecnologie necessarie per la creazione dell'applicazione mobile sono: su Android Java e la libreria android-SDK, su IOS Objective-C, su windowsPhone .NET. Tutti questi linguaggi di programmazione sono largamente documentati, elemento che faciliterebbe sicuramente il lavoro del gruppo;
	\item Il capitolato può essere molto stimolante grazie al fatto che il proponte da la possibilità al gruppo di sviluppare un'idea originale per l'applicazione. 
\end{itemize}
\subsubsection{Fattori di rischio}

Si è deciso di non accettare il capitolato per i seguenti motivi:
\begin{itemize}
\item Scarso interesse da parte dei membri del gruppo di lavoro nella teconologia TTS;
\item Difficoltà nel trovare un idea per un applicazione che fosse innovativa e utile;
\item Le Tecnologie necessarie alla realizzazione del capitolato dovrebbero essere tutte studiate da zero con il rischio di non riuscire ad avere una stima approssimativa della mole di lavoro da svolgere.
\item Mancanza di conoscenze nell'implementazione di meccanismi atti a risolvere le problematiche legate all'utilizzo di un servizio remoto.
\item L'obbiettivo di rendere disponibile la sintesi vocale a tutte le applicazioni e fornire un interfaccia di configurazione estesa comporta conoscenze specifiche e uno studio particolarmente approfondito per l'installazione e della interazione tra applicazioni e funzionalità del sistema operativo.
\end{itemize}


\end{document}
